% !TEX root = ../Main.tex
\chapter{数形结合}

数形结合就是把一个\textbf{代数表达式翻译成几何图形}，借图形的直观把答案读出来。本章看三类翻译：直接画图、化成两点距离、用数轴上的距离理解绝对值。

\section{直观画图}

\ex{求 $y=x(3-2x)$ 的最大值。}

\sol{
    \textbf{法一（画图）.} $y=-2x^2+3x$ 是开口向下的抛物线，顶点处取最大。配方 $y=-2\pr{x-\tfrac34}^2+\tfrac98$，故 $x=\tfrac34$ 时 $y_{\max}=\dfrac98$。

    \begin{center}
    \begin{tikzpicture}[>=Stealth,scale=1.0]
      \draw[->] (-0.3,0) -- (1.9,0) node[right]{$x$};
      \draw[->] (0,-0.3) -- (0,1.6) node[above]{$y$};
      \draw[thick,blue!70!black,domain=-0.05:1.55,smooth,samples=60]
          plot (\x,{(-2*\x*\x+3*\x)});
      \draw[dashed] (0.75,0) -- (0.75,1.125) node[midway]{};
      \fill (0.75,1.125) circle (1.2pt) node[above right]{\footnotesize$\pr{\tfrac34,\tfrac98}$};
    \end{tikzpicture}
    \end{center}

    \textbf{法二（基本不等式）.} 凑出和为定值：
    \[
    y=x(3-2x)=\tfrac12\cdot 2x\cdot(3-2x)\le\tfrac12\pr{\frac{2x+(3-2x)}{2}}^2=\tfrac12\cdot\frac94=\frac98,
    \]
    等号当 $2x=3-2x$，即 $x=\tfrac34$ 时取到。

    \textbf{法三（看作关于 $x$ 的二次方程有解）.} 把 $y$ 当参数，$y=x(3-2x)$ 即
    \[
    2x^2-3x+y=0.
    \]
    它要有实根 $x$，须 $\Delta=9-8y\ge 0$，即 $y\le\dfrac98$，故 $y_{\max}=\dfrac98$。
}

\section{两点距离}

平面上点 $(x,y)$ 到点 $(a,b)$ 的距离满足
\[
d^2=(x-a)^2+(y-b)^2.
\]
凡是表达式里出现"两个平方相加（开根号）"，都可以翻译成距离。

\ex{求 $y=(x-1)^2+(x+2)^2$ 的最小值。}

\sol{
    看成动点 $P(x,x)$（在直线 $y=x$ 上）到定点 $(1,-2)$ 的距离平方 $(x-1)^2+(x+2)^2$。它的最小值就是点 $(1,-2)$ 到直线 $y=x$ 的距离的平方：
    \[
    y_{\min}=\pr{\frac{\abs{1-(-2)}}{\sqrt2}}^2=\pr{\frac{3}{\sqrt2}}^2=\frac92.
    \]
}

\ex{求 $y=\sqrt{2x^2-2x+1}+\sqrt{2x^2+2x+5}$ 的最小值。}

\sol{
    配方：$2x^2-2x+1=(x-1)^2+x^2$，$2x^2+2x+5=(x+2)^2+(x-1)^2$。于是 $y=\abs{PA}+\abs{PB}$，其中 $P(x,x)$ 在直线 $y=x$ 上，$A(1,0)$、$B(-2,1)$。$A$ 在直线下方、$B$ 在上方（两点异侧），故折线拉直时最短：
    \[
    y_{\min}=\abs{AB}=\sqrt{(1+2)^2+(0-1)^2}=\sqrt{10}.
    \]
}

\ex{求 $y=\sqrt{2x^2-2x+1}-\sqrt{2x^2+2x+5}$ 的最小值。}

\sol{
    仍记 $P(x,x)$、$A(1,0)$、$B(-2,1)$，则 $y=\abs{PA}-\abs{PB}$。要 $y$ 最小，就要 $\abs{PB}$ 比 $\abs{PA}$ 超出得最多。当 $P=(1,1)$ 时 $PA$ 竖直、$PB$ 水平，$\abs{PA}=1$、$\abs{PB}=3$，取到
    \[
    y_{\min}=1-3=-2.
    \]
}

\section{绝对值}

回顾：数轴上 $\abs{x-a}$ 就是点 $x$ 到点 $a$ 的距离。求若干绝对值之和的最小值，就是在数轴上找一点，使它到各定点的距离和最小。

\begin{center}
\begin{tikzpicture}[>=Stealth,scale=1.0]
  \draw[->] (-2,0) -- (3,0) node[right]{};
  \fill (0.6,0) circle (1.2pt) node[below]{\footnotesize$a$};
  \fill (-1.2,0) circle (1.2pt) node[below]{\footnotesize$x$};
  \draw[<->,red!70!black] (-1.2,0.25) -- (0.6,0.25) node[midway,above]{\footnotesize$\abs{x-a}$};
\end{tikzpicture}
\end{center}

\ex{求 $y=\abs{x}+\abs{x-1}$ 的最小值。}

\sol{
    $\abs{x}+\abs{x-1}$ 是数轴上点 $x$ 到 $0$ 与到 $1$ 的距离之和。当 $x$ 落在 $0$ 与 $1$ 之间（含端点）时，距离和就等于两点间距 $1$，最小。
    \[
    y_{\min}=1\quad(0\le x\le 1).
    \]
}

\ex{求 $y=\abs{x-1}+\abs{x-2}+\cdots+\abs{x-9}$ 的最小值。}

\sol{
    $9$ 个定点 $1,2,\dots,9$，奇数个，取\textbf{中位数} $x=5$ 时距离和最小：
    \[
    y_{\min}=(4+3+2+1+0+1+2+3+4)=20.
    \]
}

\ex{求 $y=\abs{x-1}+\abs{x-2}+\cdots+\abs{x-10}$ 的最小值。}

\sol{
    $10$ 个定点 $1,\dots,10$，偶数个，$x$ 取中间两点之间 $[5,6]$ 的任意值时最小：
    \[
    y_{\min}=(4+3+2+1+0+1+2+3+4+5)=25\quad(5\le x\le 6).
    \]
}
